\problem{MEC-0040}{40}
(IJSO รอบที่ 1 2557) เมื่อวัตถุอยู่ห่างจากโลกมากขึ้น ความเร่งเนื่องจากแรงโน้มถ่วงของโลกจะมีขนาดน้อยลง ถ้าคำนึงถึงปัจจัยนี้ จงหาว่าวัตถุจะขึ้นไปได้ระยะสูงสุดเท่าใด ถ้าโยนวัตถุจากพื้นด้วยความเร็วต้น $u$ แล้วพบว่าวัตถุขึ้นไปสูงสุดที่ระยะหนึ่งก่อนจะตกกลับมาสู่พื้น กำหนดให้ $g=9.8\,\mathrm{m/s^2}$ ที่ผิวโลก และไม่ต้องคำนึงถึงแรงต้านอากาศ
\par\smallskip\par\smallskip\begin{tabular}{@{}p{.48\linewidth}p{.48\linewidth}@{}}1. $\dfrac{u^2}{20}$ & 2. $\dfrac{u^2}{19.6}$ \\ 3. น้อยกว่า $\dfrac{u^2}{19.6}$ & 4. มากกว่า $\dfrac{u^2}{19.6}$\end{tabular}
